\documentclass[11pt]{article}
\usepackage[margin=1in]{geometry}
\usepackage{amsmath,amssymb}
\usepackage[hidelinks]{hyperref}
\emergencystretch=3em

\newcommand{\whotimes}{\widehat{\otimes}}
\newcommand{\ellthree}{\ell_3}

\title{Literature-Implied Negative Answer to a Triple Rademacher Tensor Formula in arXiv:1103.2090}
\author{Run fa\_banach\_001, lane 19}
\date{June 15, 2026}

\begin{document}
\maketitle

\section*{Status}

This packet records a literature-implied negative subcase for a question in
Pisier, arXiv:1103.2090, \emph{Random Series of Trace Class Operators}. The
negative subcase follows from the later theorem of Briet--Naor--Regev,
arXiv:1208.0539, together with the Maurey--Pisier cotype criterion recalled in
the source paper.

\section*{Source Question}

In the final remarks of arXiv:1103.2090, after the two-fold formula
\[
R(E\whotimes F) \approx R(E)\whotimes F + E\whotimes R(F),
\]
Pisier asks whether there is a similar result for
$R(E\whotimes F\whotimes G)$. The displayed guess is
\[
R(E)\whotimes F\whotimes G
+E\whotimes R(F)\whotimes G
+E\whotimes F\whotimes R(G).
\]

\section*{Supporting Result}

Briet--Naor--Regev, Theorem 1.1, proves that if
$p_1,p_2,p_3\in(1,\infty)$ and
\[
\frac1{p_1}+\frac1{p_2}+\frac1{p_3}\leq 1,
\]
then
\[
\ell_{p_1}\whotimes \ell_{p_2}\whotimes \ell_{p_3}
\]
does not have finite cotype. In particular, with
$p_1=p_2=p_3=3$, the space
\[
B=\ellthree\whotimes \ellthree\whotimes \ellthree
\]
does not have finite cotype.

\section*{Identification}

The theorem above contradicts the guessed triple Rademacher formula in the
natural type-2 UMD class.

Indeed, $\ellthree$ has finite cotype, so the Rademacher and Gaussian sequence
spaces $R(\ellthree)$ and $G(\ellthree)$ coincide up to equivalent norms. If
Pisier's guessed formula held for
$B=\ellthree\whotimes\ellthree\whotimes\ellthree$, then every sequence in
$R(B)$ would be a sum of three projective-tensor terms, for example a term in
$R(\ellthree)\whotimes\ellthree\whotimes\ellthree$, and similarly for the
other two coordinates. Replacing $R(\ellthree)$ by $G(\ellthree)$ and using
the projective tensor norm shows that each such term is Gaussian-convergent in
$B$. Hence $R(B)\subseteq G(B)$.

The reverse inclusion $G(B)\subseteq R(B)$ holds in every Banach space. Thus
the guessed formula would imply $R(B)=G(B)$. By the Maurey--Pisier criterion
quoted in the source paper, this is equivalent to finite cotype of $B$.
This contradicts Briet--Naor--Regev's Theorem 1.1 for $p_1=p_2=p_3=3$.
Therefore the proposed triple formula cannot hold in general, even for the
classical choice $E=F=G=\ell_3$.

\section*{Scope}

This is not a new theorem; it is a direct implication from later literature.
The supporting paper explicitly studies finite-cotype questions for projective
tensor products and cites the random-series source as one place where such
questions were stated, but it does not phrase its theorem as a disproof of the
triple $R(E\whotimes F\whotimes G)$ formula.

This packet does not answer the Hilbert triple problem
$\ell_2\whotimes\ell_2\whotimes\ell_2$, the two-fold sub-2 question for
$L_p\whotimes L_q$, or the separate second-derivative formula for
$R_1(R_2(E\whotimes F))$.

\section*{References}

\begin{enumerate}
\item G. Pisier, \emph{Random Series of Trace Class Operators},
arXiv:1103.2090.
\item J. Briet, A. Naor, and O. Regev, \emph{Locally decodable codes and the
failure of cotype for projective tensor products}, arXiv:1208.0539.
\end{enumerate}

\end{document}
